\documentclass{article}

\usepackage{amsmath,amssymb}
\usepackage{xurl}
\usepackage[hidelinks]{hyperref}
\setlength{\emergencystretch}{2em}

% Shared commands for notes in docs/paper-gaps/.

\DeclareUnicodeCharacter{2080}{\ifmmode{}_{0}\else\textsubscript{0}\fi}
\DeclareUnicodeCharacter{2081}{\ifmmode{}_{1}\else\textsubscript{1}\fi}
\DeclareUnicodeCharacter{2082}{\ifmmode{}_{2}\else\textsubscript{2}\fi}
\DeclareUnicodeCharacter{2099}{\ifmmode{}_{n}\else\textsubscript{n}\fi}

\newcommand{\Lean}{\textsc{Lean}}
\ExplSyntaxOn
\NewDocumentCommand{\leanid}{m}{
  \ifmmode
    \textsf{\detokenize{#1}}
  \else
    \group_begin:
    \urlstyle{sf}
    \tl_set:Nn \l_tmpa_tl {#1}
    \tl_replace_all:Nnn \l_tmpa_tl {\_} {_}
    \exp_args:NV \path \l_tmpa_tl
    \group_end:
  \fi
}
\ExplSyntaxOff
\providecommand{\tr}{\operatorname{tr}}
\newcommand{\tnleanIssueBase}{https://github.com/LionSR/TNLean/issues/}
\newcommand{\tnleanPrBase}{https://github.com/LionSR/TNLean/pull/}
\newcommand{\tnleanDocBase}{https://sirui-lu.com/TNLean/docs/}
\newcommand{\ghissue}[1]{\href{\tnleanIssueBase#1}{GitHub issue \##1}}
\newcommand{\ghpr}[1]{\href{\tnleanPrBase#1}{PR \##1}}
\newcommand{\doclink}[1]{\href{\tnleanDocBase#1.html}{\path{#1}}}

% Dirac notation and the identity operator, as in blueprint/src/macros/common.tex.
% \providecommand keeps note-local definitions authoritative.
\usepackage{dsfont}
\providecommand{\ket}[1]{|#1\rangle}
\providecommand{\bra}[1]{\langle #1|}
\providecommand{\braket}[2]{\langle #1 | #2 \rangle}
\providecommand{\ketbra}[2]{|#1\rangle\!\langle #2|}
\providecommand{\Id}{\mathds{1}}
\providecommand{\one}{\mathds{1}}
\providecommand{\C}{\mathbb{C}}
\providecommand{\MN}[1]{M_{#1}(\C)}
\providecommand{\MD}{M_D(\C)}

% External referencing for paper sources.  The build helper writes sanitized
% auxiliary files under build/paper-aux/ and excludes duplicated source labels.
\usepackage{xr}

% Import a generated paper auxiliary file with a caller-supplied label prefix.
% The candidate paths support builds from docs/paper-gaps/, the repository root,
% and build/paper-aux/ itself.
\newcommand{\importpaperaux}[2]{%
  \IfFileExists{../../build/paper-aux/#2.aux}{%
    \externaldocument[#1]{../../build/paper-aux/#2}%
  }{%
    \IfFileExists{build/paper-aux/#2.aux}{%
      \externaldocument[#1]{build/paper-aux/#2}%
    }{%
      \IfFileExists{#2.aux}{%
        \externaldocument[#1]{#2}%
      }{}%
    }%
  }%
}

% General external reference macros
\newcommand{\paperref}[2]{\ref{#1-#2}}
\newcommand{\papereqref}[2]{(\ref{#1-#2})}

% CPSV16 source (1606.00608 / MPDO-22-12-17-2.tex) external document and shortcuts
\importpaperaux{cpsv16-}{1606.00608/MPDO-22-12-17-2-tnlean}
\newcommand{\cpsvref}[1]{\paperref{cpsv16}{#1}}
\newcommand{\cpsveqref}[1]{\papereqref{cpsv16}{#1}}

% Verdict marker: \gapnote{<kind>}{<status>}. Kind is the policy
% classification (clarification, local-correction, scope-restriction,
% unfaithful, false-source, open-gap); status is open, wip (elimination
% underway), resolved, or historical. Machine-read by the site index;
% prints a verdict line.
\newcommand{\gapnote}[2]{\par\noindent\textbf{Verdict:} #1 (#2).\par}


\title{The Abstract Schwarz-Map Hypotheses in Wolf Theorem 6.12}
\author{}
\date{2026-07-18}

\begin{document}
\maketitle

\section{Source statement}

Wolf introduces a Schwarz map as a positive unital matrix map satisfying
\begin{align}
  E(A^\dagger A)-E(A^\dagger)E(A)&\succeq0
  \qquad\text{for every }A.
  \tag{5.2}
\end{align}
This terminology is stated in Chapter~5 immediately before Example~5.3
(local source lines 347--352).  Theorem~6.12 then assumes a unital map satisfying
this Schwarz inequality and an arbitrary full-rank fixed point of its
trace-pairing adjoint.  Its algebraic conclusion is
\begin{align}
  \operatorname{Fix}(E)&=\{X:E(X)=X\}
  \quad\text{is a unital $*$-subalgebra.}
  \tag{6.12}
\end{align}
The proof is at local source lines 1395--1417.  It invokes the proposition on
positive fixed points (local lines 1161--1192, cited there as Proposition~6.9)
to replace the full-rank witness by a positive-definite fixed point before
performing the weighted-trace argument.
The complete theorem also invokes the explicit block realization in
Equation~(1.39); the present increment concerns only the displayed
$*$-algebra conclusion.

The word ``positive'' is not repeated in the theorem statement because it is
part of the Chapter~5 definition of a Schwarz map.  Exposing positivity as a
separate formal hypothesis therefore does not strengthen the source theorem.

\section{The weighted equality}

Let $\rho>0$ satisfy $E^*(\rho)=\rho$, and let $E(X)=X$.  Positivity implies
$E(X^\dagger)=X^\dagger$.  For the Schwarz gap
\begin{align}
  \Delta_X&=E(X^\dagger X)-E(X^\dagger)E(X)\succeq0,
\end{align}
the trace-pairing identity gives Wolf's Equation~(6.60):
\begin{align}
  \tr(\rho\Delta_X)
    &=\tr\!\left(E^*(\rho)X^\dagger X\right)
      -\tr\!\left(\rho X^\dagger X\right)\\
    &=0.
\end{align}
Faithfulness of the positive-definite weight forces $\Delta_X=0$.  Repeating
the calculation for $X^\dagger$ gives equality in both one-sided Schwarz
inequalities.  The abstract multiplicative-domain theorem then yields
\begin{align}
  E(XY)&=E(X)E(Y)=XY
\end{align}
for fixed points $X$ and $Y$.

This argument is formalized for arbitrary positive unital Schwarz maps once a
positive-definite invariant weight has been supplied explicitly.
\footnote{Formal declarations:
  \leanid{SchwarzMap.schwarz_equality_of_mem_fixedPoints},
  \leanid{SchwarzMap.mem_multiplicativeDomain_of_mem_fixedPoints}, and
  \leanid{SchwarzMap.fixedPointsStarSubalgebra} in
  \doclink{TNLean/Channel/FixedPoint/AbstractAlgebra}.}

\section{Comparison with the Kraus specialization}

The earlier declaration assumed a Kraus representation.
\footnote{Formal declaration:
  \leanid{Kraus.fixedPointsStarSubalgebra} in
  \doclink{TNLean/Channel/FixedPoint/Algebra}.}
That result is mathematically correct, but complete positivity is stronger than
the Schwarz-map hypothesis in Theorem~6.12.  It is retained as a useful
specialization of the post-reduction lemma and is not presented as the literal
formalization of the theorem.

\section{Verdict}

The earlier Kraus-only attribution was a \emph{scope restriction}, and the
abstract lemma removes that restriction.  A second scope restriction remains:
the formal signature assumes a positive-definite fixed point, while Wolf
Theorem~6.12 assumes only a full-rank fixed point and derives the faithful
positive witness from its proposition on positive fixed points.  Consequently,
the present result is the source-faithful post-reduction algebra lemma, not yet
the literal formalization of Theorem~6.12.

The elimination plan is to formalize the positive-fixed-parts proposition for
the abstract positive trace-preserving trace adjoint, derive a positive-definite
fixed point from the full-rank witness, and wrap the present lemma.  The
explicit trace-adjoint Equation~(1.39) realization remains open at this level
of generality.  The downstream Schrödinger-picture density-block
classification, support reduction, and zero-block assembly required for
Theorem~6.14 are complete; see
\path{docs/paper-gaps/wolf_theorem6_14_fixed_point_projection_gap.tex}.

\bibliographystyle{alpha}
\bibliography{references}

\end{document}
